\section{The integral torsion exponent for the truncated period module}

% BEGIN W81T1
\paragraph{W81T1. The prescribed integral quotient and its single graded pieces.}
Write $p$ for the local prime, $K_j=K_0(\mu_{p^j})$,
$G_j=G_{K_j}$, $H=G_{K_\infty}$ and $C=\mathbb C_p$. Use the
chosen tied Fontaine and Wach records of L78W13 on a good layer
$(K,s,L)$. Thus $K\subset K_s$, $s\ge1$ and $s\ge2$ if $p=2$.
Let $T$ be the underlying finite free $\mathbb Z_p$-lattice of
a source object satisfying $\mathrm{HypW}(M_0,h_0)$.
The integral coefficient module on a window $a<b$ is
\[
 X_{a,b}=T\otimes_{\mathbb Z_p}
         (\xi^a A_{\inf}/\xi^b A_{\inf}).
\]
Integral negative powers here mean the corresponding fractional
modules inside the large Witt ring, where $\xi$ is a unit.
Multiplication by $\xi^{-a}$ identifies their underlying quotients
with $A_{\inf}/\xi^{b-a}$. This is the actual $Q_{a,b}$ with its
transported Galois action, including the action on $\xi^a$ itself.
LS4--LS5 give $p$-torsion-freeness, $p$-adic completeness, and
after inverting $p$ the identification with the corresponding
filtered quotient $\mathrm{Fil}^aB_{\rm dR}/\mathrm{Fil}^bB_{\rm dR}$.
Indeed completion of $A_{\inf}[1/p]$ along the regular element
$\xi$ gives precisely its original quotients by $\xi^j$ for each
finite $j$. Tensoring with the finite free $T$ preserves these
statements and the exact sequences cutting a window at any integer.

For $g\in G_s$, $\theta(g\xi/\xi)=\chi(g)$. In fact
$t=\log[\epsilon]$ satisfies $gt=\chi(g)t$. Since
$[\epsilon]-1=\xi([\epsilon_1]-1)$, the residue of $t/\xi$ is
$\zeta_p-1$, which is nonzero and fixed by $g$. Reduce the identity
$(g\xi/\xi)g(t/\xi)=\chi(g)t/\xi$ in the de Rham DVR.
The ratio $g\xi/\xi$ is an integral unit since $g$ preserves the
kernel of $\theta$. Division by $\xi^i$ followed by the
surjective $\theta:A_{\inf}\to O_C$ therefore identifies
$X_{i,i+1}$ with $T\otimes O_C(i)$, integrally and equivariantly.
This also explains the positive sign of $i$ in this coefficient twist.
% END W81T1

% BEGIN W81T2
\paragraph{W81T2. Integral approximation to a primitive on one graded piece.}
Let $M(\sigma)$ be the action on the chosen $\mathbb Z_p$-basis of
$T$ of rank $d$, and put
\[
 m=2s+1+\min\{n\ge0:M_0+h_0\le p^n(p-1)\},\quad
 J=\mathrm{Jset}(T),\quad
 B=O_C^d.
\]
We use all the components of the source filtration. In particular
$J$ contains zero whenever zero is an actual jump. A79L6 constructs
on this very $\mathrm{HypW}$ lattice an invertible matrix $D$ with
$pB\subseteq DB\subseteq B$ and
$U(\sigma)=D^{-1}M(\sigma)\sigma(D)\in\mathrm{GL}_d(O_{K_m})$
on $G_m$. The representation $U$ is trivial on $H$, and there is
a $g\in G_m$ with $v_p(\chi(g)-1)=m$ and
$\|U(g)-1\|<|\zeta_p-1|$. A79L7, before its projection to one
embedding, gives a $C$-basis of vectors $b$ of weights $u\in J$
satisfying $\chi(\sigma)^uM(\sigma)\sigma(b)=b$ on $G_m$.
Here the required de Rham comparison on the given good layer
follows from its restricted crystallinity as in A79L7.
The weights lie in the finite interval $[-M_0,h_0]$.

Fix an integer $i$ and act on $C^d$ by
$\sigma*x=\chi(\sigma)^i M(\sigma)\sigma(x)$.
Suppose $x\in C^d$ has $\sigma*x-x\in B$ on $G_m$.
In $U$-coordinates let $y=pD^{-1}x$. Its differences for the
action with $U$ are integral since $pD^{-1}B\subset B$.
On $H$ this action is coordinatewise, so A79L2 gives
$py=z+b_0$ with $z\in K_\infty^d$, $b_0\in B$.
Let $A=\chi(g)^i U(g)$ and $\Gamma=g|_{\widehat K_\infty}$.
Then $(A\Gamma-1)z$ is integral. Split
$z=z_f+z_c$ in the normalized-trace splitting of A79L3, into
$K_m^d$ and the rational complementary summand. Both projections
are integral and commute with $A\Gamma$. In particular
$(A-1)z_f$ is integral and $p z_c$ is integral by the complementary
inverse of A79L3, since $p/(\zeta_p-1)$ is integral.

The vectors $D^{-1}b$ of weight $u$ for $U$ lie in
$\widehat K_\infty^d$ by $H$-invariance and A79L2. Their
complementary parts vanish by the injectivity of
$\chi(g)^uU(g)\Gamma-1$ there, from A79L3. They are thus eigenvectors
over $K_m$ for $A$ with eigenvalues $\chi(g)^{i-u}$; in particular
they span $C^d$. This argument retains the eigenspace with $u=i$.
% END W81T2

% BEGIN W81T3
\paragraph{W81T3. The resonant weight costs no denominator.}
Put
\[
 F_i(Z)=\prod_{u\in J,\ u\ne i}(Z-\chi(g)^{i-u}),\qquad
 r_i=\sum_{u\in J,\ u\ne i}(m+v_p(u-i)),\quad n_i=\max(1,r_i).
\]
The polynomial has integral coefficients and
$(A-1)F_i(A)=0$ by the basis in W81T2. The binomial valuation
calculation in A79L4 gives $v_p(F_i(1))=r_i$, with $F_i(1)=1$
when the product is empty. As $F_i(Z)-F_i(1)$ is an integral
multiple of $Z-1$, the integrality of $(A-1)z_f$ shows that
\[
 z_* = \frac{p^{n_i}}{F_i(1)}F_i(A)z_f
       \quad\text{satisfies}\quad p^{n_i}z_f-z_*\in O_{K_m}^d .
\]
We have $Az_*=z_*$. It is fixed for the twisted $U$-action by $H$
and by $g$, because its coordinates lie in $K_m$. The image of
$g$ generates a dense subgroup in $\mathrm{Gal}(K_\infty/K_m)$:
$v_p(\chi(g)-1)=m$, with $m\ge2$ at $p=2$, makes its successive
$p$-powers generate each successive quotient. Continuity therefore
gives invariance under all of $G_m$.

Now $p^{n_i}z_c\in B$, and $py=z+b_0$ in W81T2. Multiplication
back by $D$ shows
\[
                  p^{2+n_i}x-Dz_*\in B .
\]
The correction $Dz_*$ is invariant for the twisted original action.
By W81T1 this proves the following assertion about the actual
graded piece: every rational primitive with integral differences
becomes integral modulo a rational invariant after multiplication
by $p^{2+n_i}$. In particular a torsion continuous cocycle of
this piece is killed in cohomology by that power: divide the
integral primitive of a killing multiple by this multiple and use
the assertion just proved. The use of distinct weights in $F_i$
makes the exponent independent of their multiplicities.
% END W81T3

% BEGIN W81T4
\paragraph{W81T4. Lifting rational invariants without increasing the exponent.}
For $a\le j\le b$ the map $X_{a,b}[1/p]\to X_{a,j}[1/p]$
is surjective on $G_m$-invariants. To see this, write
$V=T[1/p]$ and take the adapted basis of invariant de Rham periods
$w_k$ from A79L7, of filtration degrees $u_k$. Its leading vectors
$b_k=t^{-u_k}w_k\bmod t$ form a $C$-basis as in W81T2.
The invariants in $V\otimes\mathrm{gr}^q B_{\rm dR}$
are precisely the $K_m$-span of the leading symbols of those
$w_k$ with $u_k=q$. Indeed after extracting the formal factor $t^q$
expand in the basis $b_k$. The coefficient at $k$ then satisfies
$\sigma(c)=\chi(\sigma)^{u_k-q}c$ on $G_m$. It is fixed by $H$
and hence belongs to $\widehat K_\infty$ by A79L2. The complementary
part in A79L3 is zero by its scalar twisted inverse. On $K_m$ the
generator $g$ forces $c=0$ unless $u_k=q$, because $\chi(g)$ has
infinite order. For equal weights the coefficient can be any
element of $K_m$, embedded in the de Rham period ring in the
given way.

Start with an invariant class modulo $\mathrm{Fil}^j$ at filtration
$a$. Its image in degree $a$ is invariant. Subtract the just-described
linear combination of actual invariant periods of degree $a$ to
increase the filtration by one. Repeat through degree $j-1$.
The sum of the subtracted periods gives an invariant representative
in $V\otimes\mathrm{Fil}^a B_{\rm dR}$. Its reduction modulo
$\mathrm{Fil}^b$ gives the desired lift by W81T1.
In particular an integral invariant class lifts after multiplication
by some power of $p$ to an integral invariant class. This power is
used only to identify rational invariant subspaces; it is never
added to the numerical torsion exponent.
% END W81T4

% BEGIN W81T5
\paragraph{W81T5. Adding the graded exponents on the full window.}
For a stable $p$-torsion-free module $X$ call a rational vector
$x\in X[1/p]$ integral modulo invariants at exponent $c$ if
$p^cx-x_0\in X$ for some $x_0\in X[1/p]^{G_m}$.
In an equivariant exact sequence
$0\to X'\to X\xrightarrow{q}X''\to0$ of the present modules,
suppose the assertion for every vector with integral differences
holds at exponents $c'$ and $c''$ on the two ends, and $q$ is
surjective on rational invariants. For such a vector $x$ in the
middle, we may write $p^{c''}q(x)=u''+z''$ with $u''$ invariant
and $z''$ integral. Lift $u''$ to a rational invariant $u$ and
$z''$ to an integral element $z$. Then
$y=p^{c''}x-u-z\in X'[1/p]$ has integral differences in $X'$:
they belong to $X\cap X'[1/p]=X'$, since $X''$ is torsion-free.
Apply the assertion at $c'$ to $y$ and add $p^{c'}z$ back. This
proves the middle assertion with exponent $c'+c''$.
The proof does not choose an integral lift of an invariant with a
specified bound.

Apply this argument to
$0\to X_{b-1,b}\to X_{a,b}\to X_{a,b-1}\to0$.
W81T4 gives precisely its rational invariant lifting and W81T3
is the one-step case. Thus every rational primitive with integral
differences on $G_m$ is integral modulo invariants at exponent
$C_{a,b}=\sum_{a\le i<b}(2+n_i)$.
To descend to $G_s$, use a set $R$ of representatives of the
left cosets $G_s/G_m$. For a rational primitive $x$ with integral
differences on $G_s$, choose $u\in X_{a,b}[1/p]^{G_m}$ with
$p^{C_{a,b}}x-u$ integral. The sum $\sum_{r\in R}r u$ is
$G_s$-invariant, and
$[G_s:G_m]p^{C_{a,b}}x-\sum_r r u$ is integral since each
$p^{C_{a,b}}(x-rx)$ is integral. Eisenstein's cyclotomic degree
formula over the unramified $K_0$ gives
$[G_s:G_m]=[K_m:K_s]=p^{m-s}$, including $m=s$.
The difference of the $p$-valuations of their absolute degrees
is therefore $m-s$, exactly $\mathrm{vdegMs}(m)$.
% END W81T5

% BEGIN W81T6
\paragraph{W81T6. The torsion multiplier consumed by the ramified extension.}
Fix any continuous cocycle $c$ of $G_s$ in $X_{a,b}$ for which
$p^kc$ is an integral coboundary. If $v$ is its primitive, then
$x=p^{-k}v$ has differences $c$ because $X_{a,b}$ is
$p$-torsion-free. W81T5 now gives an integral primitive of
$p^{N_{a,b}}c$ with
\[
 N_{a,b}=m-s+
       \sum_{a\le i<b}\left(2+\max\left(1,
               \sum_{\substack{u\in J\\u\ne i}}
                      (m+v_p(u-i))\right)\right).
\]
This is the source's $\mathrm{cLL}$ for all $a<b$ on the chosen
records. In particular in the extension-class argument take $a=0$,
$b=e+1$, $T$ the kernel lattice, and
$c(g)=f(g)\otimes[1]$ on $G_s$. Its assumed torsion and its
cocycle identity suffice to obtain the actual integral coboundary
in $T\otimes Q_{0,e+1}$ after multiplication by $p^{N_{0,e+1}}$.
The additional terms in $\mathrm{cTB}$ belong to the subsequent
passage from this coboundary to Bloch--Kato finiteness.
% END W81T6
